Consider two groups with identically-distributed skills
and characterized by different noise levels in screening. 
Our results demonstrate that if a regulatory body
(e.g., policymakers or a regulator) 
insists on the same number of tests 
and the same decision rule for both groups, 
this would yield higher false positive rates in any threshold policy. 
As a result, hired candidates from the noisier group 
would suffer higher rates of firing.
In turn, this might lead employers to erroneously conclude 
that this group's skill level is lower than it actually is.
This paper presents a policy that handles this problem 
by minimizing the false positive rates of both groups, 
in the form of a greedy policy. 
Moreover the greedy policy is efficient,
minimizing the expected number of tests per hire
among all policies that achieve a specified false positive rate and continue testing every candidates that appear better than the a new one.
However, the dynamic policy will still suffer 
(as does the simple threshold policy)
from higher false negative rates for the noisier group,
violating a notion of fairness dubbed 
\emph{equality of opportunity} in the recent literature 
on fairness in machine learning \cite{hardt2016equality}. 
We addressed this problem by modifying the greedy policy 
to reject candidate
% Namely, we reject a candidate 
iff 
$\Pr[y_i=+1|\hat{y}_{i,1}\dots \hat{y}_{i,\tau}]<\epsilon'$ by setting $\epsilon' < p$.
Our greedy policy can be made  forgiving and equalize false negative rates across groups.
In future work, we plan to explore extensions to the Gaussian model.